\question{}

Consider the time integration of the following diffusion equation:
\[ \frac{\partial u(\theta, t)}{\partial t} = \alpha \frac{\partial^2 u(\theta, t)}{\partial \theta^2}. \]
The solution is also assumed to be cyclic, i.e., $u(\theta + 2\pi, t) = u(\theta, t)$ $\forall \theta \in \mathbb{R}$.
The solution is discretized over $n$ equally spaced points in $\theta$:
$\theta_1 = 0, \theta_2 = 2\pi/n, \theta_3 = 4\pi/n, \ldots$.

You have written the following OpenCL kernel to integrate
the equation over \lstinline|num_iter| time steps.
The constant $F$ is $\alpha \Delta t / (\Delta x)^2$ where $\Delta x = 2\pi/n$.

\textbf{Simplifying assumption} \emph{This kernel assumes that the number of processing elements of a compute unit is at least $n$
so all threads are in the same compute unit.
You do not have to consider any situation in which this is not the case.}

\lstinputlisting[language=c]{Q3/ode.cl}

You observe that your kernel does not always provide the same answer,
the answer even seems to be incorrect.

Answer the following questions:
\begin{parts}
  \part \label{issue} What is the issue causing the kernel to produce incorrect results ?

    \begin{solutionbox}{9cm}
      We are missing synchronization. So a thread could be using the wrong value of \lstinline|a[prev]|
      or \lstinline|a[next]| in case they are not synchronized.

      \textbf{Note}: In view of the simplifying assumption, the local and global sizes are the same so
      \lstinline|get_local_id(0)| and 
      \lstinline|get_global_id(0)| are the same. This is not the source of issues.
    \end{solutionbox}
  \part You identify that depending on your device, if $n$ is sufficienty small
    you always get the correct answer. More precisely, you get the correct answer
    in the following situations:
    \begin{itemize}
      \item $n \le 8$ and you use your CPU as device for your kernel,
      \item $n \le 32$ and you use your GPU as device for your kernel.
    \end{itemize}
    Explain why.

    \begin{solutionbox}{9cm}
    If $n$ is that small then the execution can be run using SIMT in which case
    the threads are guaranteed to be executed each instruction at the same time.
    \end{solutionbox}
  \part How can you fix this issue you identified in \ref{issue}. ?

    \begin{solutionbox}{9cm}
      We should add \lstinline|volatile| and also add barriers to synchronize the threads.
    \end{solutionbox}
\end{parts}
